\documentclass[aspectratio=169]{beamer}

\usetheme{Madrid}
\usecolortheme{default}

\usepackage{booktabs}
\usepackage{array}
\usepackage{graphicx}

\title{Synthetic Clinical Note Generation for GLP-1 RA Discontinuation Reason Classification}
\author{Adiba, Shahidul \\ BioNLP Research Group}
\date{\today}

\begin{document}

\begin{frame}
    \titlepage
\end{frame}

\begin{frame}{Presentation Overview}
    \tableofcontents
\end{frame}

\section{Motivation and Scope}

\begin{frame}{Problem and Motivation}
    \begin{itemize}
        \item GLP-1 receptor agonists are widely used for diabetes and weight management, but discontinuation is common and clinically important.
        \item Reasons for discontinuation are often buried in free-text clinical notes rather than structured fields.
        \item Supervised NLP models need large, high-quality labeled datasets, but annotated clinical notes are expensive to produce.
        \item This project uses synthetic note generation to build training data for discontinuation reason classification.
    \end{itemize}
\end{frame}

\begin{frame}{Study Objective}
    \begin{itemize}
        \item Build a publishable end-to-end pipeline for synthetic GLP-1 discontinuation note generation, validation, and benchmarking.
        \item Create a diverse synthetic dataset in VA-style EHR format.
        \item Validate quality with embeddings, LLM judges, and human review.
        \item Benchmark the dataset on downstream classification with strong model baselines.
        \item Report rigorous dataset statistics, ablations, and limitations suitable for conference or journal submission.
    \end{itemize}
\end{frame}

\begin{frame}{Current Taxonomy and Expansion Plan}
    \textbf{Current four-category setup}
    \begin{enumerate}
        \item \textbf{Adverse Effects / Tolerability}
        \item \textbf{Access / Cost / System Barriers}
        \item \textbf{Patient Preference / Beliefs}
        \item \textbf{Planned / Goal-Based Discontinuation}
    \end{enumerate}

    \vspace{0.5em}
    \textbf{Additional categories available for future expansion}
    \begin{itemize}
        \item Clinical Decision / Medical Appropriateness
        \item Efficacy / Expectation Mismatch
        \item Other / Unknown
    \end{itemize}

    \begin{block}{Note}
        The current results use the four broad categories above; the extra seed categories support future taxonomy expansion.
    \end{block}
\end{frame}

\begin{frame}{Research Questions and Planned Contributions}
    \small
    \textbf{Research questions}
    \begin{itemize}
        \item Can synthetic GLP-1 discontinuation notes achieve strong realism, diversity, and label fidelity?
        \item How well do LLM-based judgments align with expert human review?
        \item Does the curated synthetic dataset improve downstream classification in zero-shot, few-shot, and fine-tuned settings?
        \item Which curation steps matter most for downstream performance?
    \end{itemize}

    \textbf{Expected contributions}
    \begin{itemize}
        \item A new clinically motivated synthetic dataset and benchmark
        \item A validation pipeline combining embeddings, LLM judges, and human judges
        \item A multi-model benchmark for GLP-1 discontinuation reason classification
        \item An empirical study of when synthetic data helps and where it fails
    \end{itemize}
\end{frame}

\section{Methodology and Study Design}

\begin{frame}{Synthetic Data Pipeline}
    \begin{enumerate}
        \item \textbf{Generation:} Few-shot prompting of LLMs to create VA-style clinical notes
        \item \textbf{Diversity control:} Qwen3-Embedding-8B deduplication with cosine similarity thresholding
        \item \textbf{Automated validation:} LLM-as-a-Judge scoring on five quality dimensions
        \item \textbf{Human validation:} Expert review on a sampled subset for quality and label correctness
        \item \textbf{Curation:} Threshold-based filtering plus creation of train/dev/test partitions
    \end{enumerate}

    \vspace{0.5em}
    \begin{block}{Generation setup}
        Temperature 0.9, top-p 0.95, max output 4,192 tokens, max model length 8,192, seed 42.
    \end{block}
\end{frame}

\begin{frame}{Evaluation Plan}
    \small
    \begin{itemize}
        \item Generate and curate the synthetic dataset
        \item Compute dataset statistics: class balance, note length, lexical diversity, embedding diversity, and judge score distributions
        \item Reserve \textbf{200 notes} as a held-out test set and keep them untouched during development
        \item Measure agreement between LLM-as-a-Judge and human-as-a-Judge review
        \item Use ablations to test the value of diversity filtering, LLM filtering, and human filtering
    \end{itemize}
\end{frame}

\begin{frame}{Benchmark Models and Settings}
    \small
    \textbf{Planned model families}
    \begin{itemize}
        \item Qwen3-4B
        \item Qwen3-8B
        \item Qwen3.5-9B
        \item Gemma 4
        \item Llama 3.1
        \item Llama 3.2
        \item Encoder baselines such as BioClinicalBERT or ClinicalBERT
    \end{itemize}

    \textbf{Planned settings}
    \begin{itemize}
        \item Zero-shot
        \item Few-shot
        \item Supervised / fine-tuned
        \item DPO or other preference-based optimization
    \end{itemize}
\end{frame}

\begin{frame}{Publication-Ready Study Requirements}
    \small
    \begin{itemize}
        \item \textbf{Human evaluation protocol:} reviewer expertise, annotation rubric, sample size, and agreement metrics such as Cohen's kappa or Krippendorff's alpha
        \item \textbf{Strong baselines:} include encoder baselines such as BioClinicalBERT or ClinicalBERT, not only generative LLMs
        \item \textbf{Ablation studies:} test the effect of diversity filtering, LLM filtering, human filtering, prompt design, and category granularity
        \item \textbf{Agreement analysis:} measure alignment between LLM-as-a-Judge and human-as-a-Judge decisions
        \item \textbf{External validation:} compare against real or expert-authored notes when possible
        \item \textbf{Release standards:} dataset card, reproducibility details, privacy/ethics checks, bias analysis, and strict test-set discipline
    \end{itemize}
\end{frame}

\section{Current Results}

\begin{frame}{Generation Runs}
    \small
    \begin{table}
        \centering
        \begin{tabular}{llll}
            \toprule
            \textbf{Model} & \textbf{Run} & \textbf{Notes/Category} & \textbf{Total} \\
            \midrule
            Gemma-4-31B-it & v1 exploratory & 100 & 400 \\
            Qwen3.5-27B & v1 exploratory & 100 & 400 \\
            Qwen3.6-27B & v1 exploratory & 81--100 & 280 \\
            GLM-4.7-Flash & v1 exploratory & 100 & 400 \\
            \midrule
            \multicolumn{3}{l}{\textbf{Subtotal exploratory}} & \textbf{1,480} \\
            \midrule
            Gemma-4-31B-it & v2 scale-up & 1,000 & 4,000 \\
            \midrule
            \textbf{Grand Total} & & & \textbf{5,480} \\
            \bottomrule
        \end{tabular}
    \end{table}

    \begin{itemize}
        \item The exploratory phase compared four generation models.
        \item Gemma-4-31B-it was selected for scale-up after the strongest quality results.
    \end{itemize}
\end{frame}

\begin{frame}{Gemma v2 Scale-Up}
    \small
    \begin{table}
        \centering
        \begin{tabular}{lrr}
            \toprule
            \textbf{Category} & \textbf{Notes} & \textbf{Elapsed Time (s)} \\
            \midrule
            Adverse Effects / Tolerability & 1,000 & 1,321.5 \\
            Access / Cost / System Barriers & 1,000 & 1,217.0 \\
            Patient Preference / Beliefs & 1,000 & 1,363.9 \\
            Planned / Goal-Based Discontinuation & 1,000 & 1,002.6 \\
            \midrule
            \textbf{Total} & \textbf{4,000} & \textbf{4,905.0} \\
            \bottomrule
        \end{tabular}
    \end{table}

    \begin{itemize}
        \item Used the same prompting and inference settings as the exploratory Gemma run.
        \item Total generation time was about 82 minutes.
        \item Combined Gemma pool after v1 + v2: 4,400 notes.
    \end{itemize}
\end{frame}

\begin{frame}{LLM-as-a-Judge Framework}
    \begin{itemize}
        \item Judge model: \textbf{Gemma-4-31B-it}
        \item Each note was scored from 1 to 5 on:
        \begin{itemize}
            \item Clinical realism
            \item Category alignment
            \item Clinical coherence
            \item Note variability
            \item GLP-1 relevance
        \end{itemize}
        \item Each note also received an overall accept/reject decision.
    \end{itemize}

    \begin{block}{Filtering criteria}
        Retain only notes with accept decisions, strong per-dimension scores, and average score $\geq 3.0$.
    \end{block}
\end{frame}

\begin{frame}{Overall Quality by Model}
    \scriptsize
    \begin{table}
        \centering
        \begin{tabular}{lccccc}
            \toprule
            \textbf{Model} & \textbf{Realism} & \textbf{Alignment} & \textbf{Coherence} & \textbf{Variability} & \textbf{Relevance} \\
            \midrule
            Gemma-4-31B-it & 4.98 & 5.00 & 5.00 & 4.89 & 5.00 \\
            Qwen3.5-27B & 5.00 & 5.00 & 5.00 & 4.40 & 5.00 \\
            GLM-4.7-Flash & 3.82 & 4.32 & 4.01 & 3.61 & 4.32 \\
            Qwen3.6-27B & 1.16 & 1.16 & 1.16 & 1.15 & 1.16 \\
            \bottomrule
        \end{tabular}
    \end{table}

    \begin{itemize}
        \item Gemma delivered the strongest and most consistent quality.
        \item Qwen3.5 remained strong overall but was weaker on note variability.
        \item GLM and Qwen3.6 showed substantial quality degradation.
    \end{itemize}
\end{frame}

\begin{frame}{Filtering Outcome}
    \small
    \begin{table}
        \centering
        \begin{tabular}{lcccc}
            \toprule
            \textbf{Model} & \textbf{Total} & \textbf{Accept Rate} & \textbf{Kept} & \textbf{Keep Rate} \\
            \midrule
            Gemma-4-31B-it & 400 & 100.0\% & 354 & 88.5\% \\
            Qwen3.5-27B & 400 & 100.0\% & 159 & 39.8\% \\
            GLM-4.7-Flash & 400 & 82.5\% & 42 & 10.5\% \\
            Qwen3.6-27B & 280 & 3.9\% & 10 & 3.6\% \\
            \midrule
            \textbf{Total} & \textbf{1,480} & \textbf{---} & \textbf{565} & \textbf{38.2\%} \\
            \bottomrule
        \end{tabular}
    \end{table}

    \begin{block}{Key takeaway}
        Gemma was the clear winner: 4.97 mean overall score and 88.5\% retention under strict filtering.
    \end{block}
\end{frame}

\begin{frame}{Category-Level Findings}
    \small
    \begin{table}
        \centering
        \begin{tabular}{lcccc}
            \toprule
             & \textbf{Gemma} & \textbf{Qwen3.5} & \textbf{GLM} & \textbf{Qwen3.6} \\
            \midrule
            Adverse Effects & 100 & 28 & 4 & 0 \\
            Access / Cost & 100 & 22 & 11 & 0 \\
            Patient Preference & 100 & 96 & 27 & 1 \\
            Goal-Based & 54 & 13 & 0 & 9 \\
            \midrule
            \textbf{Total} & \textbf{354} & \textbf{159} & \textbf{42} & \textbf{10} \\
            \bottomrule
        \end{tabular}
    \end{table}

    \begin{itemize}
        \item Patient Preference was especially strong for Qwen3.5 and GLM.
        \item Planned / Goal-Based discontinuation was the hardest category for multiple models.
        \item The strongest current training pool comes from Gemma plus selected high-quality notes from other models.
    \end{itemize}
\end{frame}

\begin{frame}{Key Findings So Far}
    \begin{itemize}
        \item Gemma-4-31B-it is the strongest generator under the current pipeline.
        \item The 565 filtered notes provide a strong starting point for classifier training.
        \item The 4,400-note Gemma pool is the main resource for next-stage curation.
        \item \textbf{Low note variability} was the most common rejection reason across models.
        \item \textbf{Clinical realism failures} were prominent for GLM-4.7-Flash and Qwen3.6-27B.
        \item \textbf{Category misalignment} was severe for Qwen3.6-27B, indicating likely prompt-following or formatting issues.
        \item Poor performance for Qwen3.6-27B appears more consistent with an implementation problem than a true model ceiling.
    \end{itemize}
\end{frame}

\begin{frame}{Current Limitations}
    \begin{itemize}
        \item Synthetic notes may still miss the messiness of real EHR text.
        \item Evaluation relied on a single judge model.
        \item Real discontinuation reasons can be multi-factorial and harder to label cleanly than the current taxonomy.
        \item Human validation, external validation, and downstream benchmarking are still in progress.
    \end{itemize}
\end{frame}

\section{Next Steps}

\begin{frame}{Next Steps Toward a Publishable Dataset Paper}
    \begin{itemize}
        \item Evaluate the 4,000-note Gemma v2 corpus with the same automated and human validation pipeline.
        \item Apply diversity filtering, quality filtering, and error analysis to the expanded corpus.
        \item Construct stable train/dev/test splits, including a fixed 200-note held-out test set.
        \item Benchmark classifier performance across Qwen3-4B, Qwen3-8B, Qwen3.5-9B, Gemma 4, Llama 3.1, and Llama 3.2.
        \item Compare zero-shot, few-shot, supervised fine-tuning, and DPO-style preference optimization.
        \item Report publishable dataset statistics, ablations, and agreement analysis between LLM judges and human reviewers.
        \item Expand beyond the current four-class setup using additional seed categories: Clinical Decision / Medical Appropriateness, Efficacy / Expectation Mismatch, and Other / Unknown.
        \item Compare against real-world EHR sources for external validity.
    \end{itemize}
\end{frame}

\begin{frame}{Conclusion}
    \begin{itemize}
        \item Total synthetic corpus generated: \textbf{5,480 notes}
        \item Exploratory corpus retained after filtering: \textbf{565 notes}
        \item Best generation model: \textbf{Gemma-4-31B-it}
        \item Gemma scale-up corpus available for next-stage curation: \textbf{4,400 notes}
    \end{itemize}

    \vspace{0.75em}
    \begin{block}{Bottom line}
        The project has a clear presentation path: motivate the task, explain the curation pipeline, show current generation results, and end with the validation and benchmarking plan needed for a strong dataset paper.
    \end{block}
\end{frame}

\appendix
\section{Category Examples}

\begin{frame}[fragile,allowframebreaks]{Example Note: Adverse Effects / Tolerability}
\tiny
\textbf{Discontinuation cue:}
\textcolor{red}{persistent GI side effects unresponsive to mitigation, anhedonia of unclear etiology, and distressing hair loss are driving discontinuation.}

\vspace{0.5em}
\begin{verbatim}
OUTPATIENT VISIT NOTE
Date of Service: 10/14/2024
Provider: Whitfield, Dana A., MD
Location: Internal Medicine - Weight Management Clinic
Patient: [REDACTED] | 52 y/o Female
CC: Follow-up, obesity management / tirzepatide taper

HPI:
Pt is a 52 y/o female with hx of obesity presenting for scheduled f/u re:
tirzepatide (Zepbound) therapy. She has been on tirzepatide for approximately
22 months, most recently at 7.5mg SQ weekly (highest available dose). Current
weight 173 lbs, down from ~220 lbs at start. States she is approximately 8 lbs
from her personal goal weight. Has maintained weight training and walking
throughout tx; recently added jogging 1-2x/week.

Pt reports a constellation of GI side effects that have been ongoing and now
significantly affecting QOL: constipation cycling, bloating, and gas at minimum
once weekly despite aggressive dietary fiber intake and hydration strategies.
Also endorses new-onset anhedonia over the past several months, which she
attributes at least partially to the medication. Additionally reports noticeable
hair thinning - estimates significant loss over treatment course, finds this
distressing. Denies nausea or vomiting. No chest pain, no dyspnea.

Pt has self-initiated a taper - split her last 7.5mg dose to 3.75mg and has
approximately 3-4 doses remaining. Presents today to discuss whether to
complete taper or stop now, and what to expect afterward. Not requesting
restart at this time.

PMH: Obesity (class II at baseline), dyslipidemia, iron deficiency (resolved)
Medications: Tirzepatide 7.5mg SQ weekly (tapering), creatine monohydrate
(OTC), multivitamin
Allergies: NKDA

Vitals: Wt 173 lbs | BMI 28.1 | BP 118/74 | HR 72 | SpO2 98%
Objective: Well-appearing female, NAD. No thyromegaly. Diffuse hair thinning
noted on exam, consistent with telogen effluvium pattern. Abdomen soft, NTND.

Labs (09/28/2024): Ferritin 22 (low-nl), TSH 1.9, CBC WNL, CMP WNL.

Assessment and Plan:
Pt near goal weight after ~22 months of tirzepatide with multiple converging
factors driving discontinuation: persistent GI side effects unresponsive to
mitigation, anhedonia of unclear etiology, distressing hair loss, and formulary
exclusion making continued access financially unsustainable. Patient has
independently initiated a dose taper, which is clinically reasonable given
duration of use.

Endorsed completing taper at 3.75mg for remaining doses rather than abrupt
stop. Counseled re: expected weight stabilization or modest regain
post-discontinuation and importance of maintaining current exercise regimen.
No new Rx today. RTC 3 months for weight check, metabolic panel, mood
reassessment.
\end{verbatim}
\end{frame}

\begin{frame}[fragile,allowframebreaks]{Example Note: Access / Cost / System Barriers}
\tiny
\textbf{Discontinuation cue:}
\textcolor{red}{insurance has removed tirzepatide from formulary and the patient cannot sustain out-of-pocket cost, so discontinuation is driven by access barriers.}

\vspace{0.5em}
\begin{verbatim}
OUTPATIENT VISIT NOTE
Date of Service: 10/14/2024
Provider: Whitfield, Dana A., MD
Location: Internal Medicine - Weight Management Clinic
Patient: [REDACTED] | 52 y/o Female
CC: Follow-up, obesity management / tirzepatide taper

HPI:
Pt is a 52 y/o female with hx of obesity presenting for scheduled f/u re:
tirzepatide (Zepbound) therapy. She has been on tirzepatide for approximately
22 months, most recently at 7.5mg SQ weekly. Current weight 173 lbs, down from
~220 lbs at start. States she is approximately 8 lbs from her personal goal
weight.

Separately, pt notes insurance has removed tirzepatide from formulary
entirely. Her plan has also moved semaglutide (Wegovy) to Tier 3, making OOP
cost comparable to full self-pay. States she cannot sustain either medication
financially at current pricing and is not pursuing self-pay.

Pt has self-initiated a taper - split her last 7.5mg dose to 3.75mg and has
approximately 3-4 doses remaining. Presents today to discuss whether to
complete taper or stop now, and what to expect afterward. Not requesting
restart at this time.

PMH: Obesity (class II at baseline), dyslipidemia
Medications: Tirzepatide 7.5mg SQ weekly (tapering), multivitamin
Allergies: NKDA
Vitals: Wt 173 lbs | BMI 28.1 | BP 118/74 | HR 72

Assessment and Plan:
Pt near goal weight after ~22 months of tirzepatide. A major driver of
discontinuation is formulary exclusion making continued access financially
unsustainable. Patient has independently initiated a dose taper, which is
clinically reasonable given duration of use.

Plan is to complete taper at 3.75mg for remaining doses rather than abrupt
stop. Reviewed expected weight stabilization or modest regain and discussed
continuing exercise and nutrition habits after discontinuation. No new Rx
today. RTC 3 months for follow-up.
\end{verbatim}
\end{frame}

\begin{frame}[fragile,allowframebreaks]{Example Note: Patient Preference / Beliefs}
\tiny
\textbf{Discontinuation cue:}
\textcolor{red}{the patient wants to stop based on online misinformation and personal beliefs, despite no physical adverse effects.}

\vspace{0.5em}
\begin{verbatim}
DATE: 01/27/2025
AUTHOR: MBEKI, FAITH N., NP
TITLE: TELEPHONE ENCOUNTER - MEDICATION MANAGEMENT

ACTIVE PROBLEMS: T2DM, anxiety (GAD), hypothyroidism, obesity,
hyperlipidemia, migraines

VITALS: Not obtained (telephone visit)

SUBJECTIVE:
Telephone encounter w/ 55 y/o female veteran. Pt calling with questions about
her semaglutide (Ozempic) 0.5mg weekly (started 8 weeks ago). States she
researched online and "read that Ozempic causes muscle loss." Concerned about
sarcopenia. Also reports her sister told her the medication is "only for rich
people and it's just hype." Pt states "I don't really feel any different on it
and I'm not sure it's doing anything." Denies nausea, vomiting, diarrhea.
Denies any physical side effects. Weight per home scale: unchanged from start.
Pt asks if she can "just take a break" from it. When asked to elaborate,
states she wants to stop for now and "maybe try again later." When probed about
specific concerns, repeats the information from online and her sister's
comments. Does not report cost issues. Does not report any physical symptoms.

OBJECTIVE: N/A - telephone encounter.
LABS: Most recent A1c 8.4% (11/2024 at initiation). No interval labs since.
ACTIVE MEDS: Semaglutide 0.5mg SQ weekly (8 wks), levothyroxine 75mcg QAM,
metformin 500mg BID, escitalopram 10mg QD, topiramate 50mg QD, atorvastatin
20mg QHS.

ASSESSMENT AND PLAN:
1. T2DM / SEMAGLUTIDE - Pt requesting to stop semaglutide. No adverse effects
reported. Pt's stated concerns appear based on online misinformation and social
influence rather than clinical experience with the drug. Counseled pt at length
re: evidence base for semaglutide and addressed muscle loss concern. Pt
listened but maintains she wants to stop "for now." Her preference documented.

- Semaglutide 0.5mg: per patient preference, will hold and not refill at this
  time. Pt understands A1c may worsen. Advised to check BG more frequently.
- Continue metformin 500mg BID.
- Repeat A1c and BMP at 3-month follow-up visit.
\end{verbatim}
\end{frame}

\begin{frame}[fragile,allowframebreaks]{Example Note: Planned / Goal-Based Discontinuation}
\tiny
\textbf{Discontinuation cue:}
\textcolor{red}{the patient has nearly reached the target weight and is ending therapy intentionally as a planned, time-limited course.}

\vspace{0.5em}
\begin{verbatim}
Ambulatory Care Visit - Weight Management Follow-up
Date: 03/06/2025
Clinician: Osei, Priscilla K., PA-C
Clinic: Adult Primary Care / Metabolic Health

34 y/o female here today for what appears to be her final scheduled visit in
our weight management program. Pt initiated semaglutide (Wegovy)
approximately 8 months ago, self-pay, with an explicitly stated short-term goal
of returning to her pre-pregnancy weight. She was counseled at initiation that
long-term use was not her plan and that she understood the medication as a
time-limited tool to help reestablish healthy habits and reach a defined target.

At today's visit she is within 5 lbs of her stated goal. Over the past several
weeks she has been deliberately extending injection intervals - moved from every
7 days to every 10 days, then to every 14 days for her most recent dose.
Reports she has continued to lose weight during this period and is not
experiencing rebound hunger or significant cravings. She attributes this to
dietary changes she has sustained throughout - structured meals, reduced
processed food intake, consistent physical activity. She does not report GI
side effects or any other complaints related to the medication.

Pt is self-pay and states the cost was always a planned, time-limited expense.
She is not requesting a new supply and does not wish to continue. No distress
about stopping. She seems well-prepared for the transition off medication.

Wt: 138 lbs (starting wt documented as ~168 lbs). BMI 22.9. BP 112/70, HR 68.
No new labs indicated at this visit.

Plan:
Pt has completed her intended course of therapy in a structured, intentional
manner. Taper approach (interval extension rather than abrupt stop) is
reasonable and she has tolerated it well. No clinical concerns at this time.
Discussed weight maintenance strategies and that some modest regain is common
post-discontinuation and does not represent failure.
\end{verbatim}
\end{frame}

\section{Synthetic Dataset Examples}

\begin{frame}[fragile,allowframebreaks]{Synthetic Example: Adverse Effects / Tolerability}
\tiny
\textbf{Discontinuation cue:}
\textcolor{red}{severe nausea, early satiety, bloating, and intermittent vomiting are significantly impairing quality of life, so semaglutide is discontinued immediately.}

\vspace{0.5em}
\begin{verbatim}
Source: Filtered synthetic note from Gemma-4-31B-it

VA PRIMARY CARE CLINIC NOTE
Date of Service: 11/12/2024
Provider: Sterling, Marcus J, MD
Location: Community Health Outpatient Clinic
Patient: [REDACTED] | 58 y/o Male

CC: Follow-up T2DM / GLP-1 tolerability

HPI:
Pt is a 58 y/o male with a hx of Type 2 Diabetes Mellitus, hypertension, and
hyperlipidemia presenting for scheduled f/u. He has been on semaglutide
(Ozempic) for approximately 6 months, currently at 0.5mg SQ weekly. HbA1c has
improved from 8.4% to 7.2% since initiation.

Pt reports significant and worsening GI distress over the last 4-6 weeks.
Specifically describes severe nausea that is now present daily, often
persisting for several hours after meals. He endorses an "unbearable" feeling
of early satiety and bloating, noting he can barely finish a small portion of
dinner without feeling ill. Reports intermittent vomiting (1-2x per week),
though denies any hematemesis or fever. He notes these symptoms are
significantly impacting his ability to work and his overall QOL.

PMH: T2DM, HTN, HL, OSA
Medications:
- Semaglutide 0.5mg SQ weekly
- Lisinopril 20mg PO daily
- Atorvastatin 40mg PO daily
- Metformin 1000mg PO BID

Vitals: Wt 218 lbs | BMI 31.4 | BP 132/84 | HR 78 | SpO2 96% on RA

Assessment and Plan:
Semaglutide intolerance with severe nausea, early satiety, and intermittent
vomiting despite dietary modifications. Symptoms are now severely impacting
quality of life. Plan is to discontinue semaglutide immediately, continue
Metformin, and monitor glycemic control closely after cessation.
\end{verbatim}
\end{frame}

\begin{frame}[fragile,allowframebreaks]{Synthetic Example: Access / Cost / System Barriers}
\tiny
\textbf{Discontinuation cue:}
\textcolor{red}{semaglutide became non-covered, the appeal was denied, and the out-of-pocket cost is unsustainable, forcing discontinuation.}

\vspace{0.5em}
\begin{verbatim}
Source: Filtered synthetic note from Gemma-4-31B-it

VA MEDICAL CENTER
Department of Primary Care / Internal Medicine
Progress Note

Patient: [REDACTED]
DOB: [REDACTED] (58M)
Date of Visit: 11/12/2024
Provider: Henderson, Marcus, MD

Reason for Visit: Follow-up for Type 2 Diabetes Mellitus and medication review.

S:
Pt is a 58 y/o male with T2DM, hypertension, and hyperlipidemia. He reports
that semaglutide (Ozempic) worked well over the last 14 months with good
glycemic control and modest weight loss. However, his insurance provider has
changed its formulary and semaglutide is now non-covered. Appeal was denied.

Other GLP-1 agonists are similarly not covered. The out-of-pocket cost for a
monthly supply is now approximately $900, which he states is completely
unsustainable. He has been stretching his last few doses while trying to find a
solution, but is now effectively out of medication.

O:
Wt 224 lbs (up 4 lbs from last visit). BMI 31.2. BP 132/84. HR 78.
Recent Labs: HbA1c 6.8%, Fasting Glucose 122, Creatinine 1.1, LDL 82.

A/P:
Pt achieved excellent glycemic control on semaglutide but now faces abrupt loss
of access due to insurance formulary exclusion. The financial barrier is
absolute. Plan is to discontinue semaglutide, increase Metformin to 1000mg BID,
and monitor closely for glycemic rebound after GLP-1 discontinuation.
\end{verbatim}
\end{frame}

\begin{frame}[fragile,allowframebreaks]{Synthetic Example: Patient Preference / Beliefs}
\tiny
\textbf{Discontinuation cue:}
\textcolor{red}{the patient wants to stop because he believes the medication is "artificial" and does not want to depend on it, despite clinical benefit.}

\vspace{0.5em}
\begin{verbatim}
Source: Filtered synthetic note from Gemma-4-31B-it

Veterans Affairs Medical Center
Department of Internal Medicine
Progress Note - Established Patient

Patient: [REDACTED]
DOB: [REDACTED] (62M)
Date: 02/14/2025
Provider: Sarah Jenkins, MD
Visit Type: Follow-up, Type 2 Diabetes Management

S:
Patient is a 62-year-old male veteran presenting for routine DM follow-up. He
has been on tirzepatide (Mounjaro) 5mg weekly for approximately 4 months. He
reports a good initial response with a weight loss of 12 lbs and improved
fasting blood glucose levels. However, he comes in today requesting to
discontinue the medication.

Patient explains that he has become increasingly concerned about the "artificial"
nature of the medication. He states he read online that GLP-1s "trick the brain"
and lead to loss of natural willpower over time. He says he does not want to
become dependent on a chemical to make healthy choices. He denies any adverse
physical side effects and understands that his blood sugars may worsen.

O:
Wt 218 lbs (down from 230 lbs). BP 138/86. HR 72.
Labs: HbA1c 7.2% (down from 8.1%), Fasting Glucose 134.

A/P:
Patient requests discontinuation based on personal beliefs regarding dependence
and the "artificial" nature of appetite suppression. Counseling provided
regarding expected metabolic benefits and risks of stopping, but the patient
remains firm in his decision. Plan is to discontinue tirzepatide and continue
Metformin with close glucose monitoring.
\end{verbatim}
\end{frame}

\begin{frame}[fragile,allowframebreaks]{Synthetic Example: Planned / Goal-Based Discontinuation}
\tiny
\textbf{Discontinuation cue:}
\textcolor{red}{the patient reached the pre-defined weight goal, tapered successfully, and is now ready for full planned discontinuation.}

\vspace{0.5em}
\begin{verbatim}
Source: Filtered synthetic note from Gemma-4-31B-it

VA Medical Center - Metabolic Health / Endocrinology
Date: 10/14/2024
Clinician: Henderson, Marcus L., MD

Patient: 52 y/o male with Type 2 Diabetes, Hypertension, and Obesity

Reason for Visit: Planned therapy transition and weight management follow-up.

HPI:
Pt began tirzepatide (Zepbound) approximately 11 months ago with a clearly
defined, short-term goal of reaching a specific weight target (about 200 lbs)
to reduce joint pain and improve mobility. Long-term medication use was never
his intended plan; the therapy was framed as a time-limited tool to facilitate
weight loss and establish sustainable diet and exercise habits.

At today's visit, Pt has reached his stated goal (Current Wt: 198 lbs). Over
the last 6 weeks, he has been intentionally tapering by extending the injection
interval from every 7 days to every 10, then 14, and most recently 21 days.
During this period he has maintained his weight and denies rebound hunger or a
return of obesity-related joint pain.

Assessment/Plan:
Pt has successfully reached his pre-defined weight goal in a structured and
intentional manner. The taper approach via interval extension has been well
tolerated, and he is now clinically ready for full discontinuation of
tirzepatide. No further prescriptions will be issued. He will continue current
diet and exercise habits and transition back to routine primary care follow-up.
\end{verbatim}
\end{frame}

\end{document}
